\section{Freedom from the Fear of Death}

The Letter to the Hebrews teaches that the Son saves us by sharing our humanity. Those whom he sanctifies are his brothers; he is not ashamed to call them so. Because the children share flesh and blood, he has shared the same human condition, entering death to break the devil's power and release those whose fear of death has held them in lifelong servitude (Heb 2:11--18). He takes the flesh in which we suffer and offers the obedience we could not supply for ourselves. His victory concerns us because he has made himself our brother and undergone the death we fear.

Earlier, Hebrews distinguishes what we see around us from what has been accomplished in Christ. We do not yet see all things subjected to humanity, but we see Jesus, who suffered death and is crowned with glory and honor (Heb 2:8--10). Violence continues, bodies fail, and the faithful remain exposed to loss. We take courage from the Lord's exaltation while awaiting the completion of his victory in his members. The suffering still before us must be endured with his help, in hope of the glory he has already attained.

Fear enslaves us when it leads us to act against conscience. Bodily survival is a great created good, and the care of life ordinarily serves love. Yet a person can begin to treat continued safety as a good for which every other good must be surrendered. We may conceal the truth because speaking it carries a cost, refuse another person's need because helping would disturb our security, or practice the faith only while it leaves our arrangements untouched. Such choices train us to avoid the sacrifices fidelity requires. The fear of death can influence our conduct long before our life is threatened, through a growing refusal to risk anything that sustains its comfort or security.

Christ addresses that captivity when he sends the disciples to bear witness. He teaches them whom to fear and what to trust: those who can kill the body cannot determine the soul's destiny; the Father knows the sparrows and has numbered the hairs of their heads (Mt 10:26--33). Divine care and the possibility of persecution appear together in his teaching. The disciples are precious to the Father within the danger of their mission. They can confess the Son because they are already known and loved by the Father, whose knowledge reaches beyond every power that may threaten them.

This confidence helps us act courageously in ordinary duties. We may be asked to speak an unwelcome truth, protect someone with little influence, or bear a loss that honest conduct makes unavoidable. The immediate calculation can be discouraging. A useful friendship may cool; advancement may become less likely; someone whose approval we value may misunderstand us. Remembering death helps us consider how much these losses should weigh against our duty to God. The people whose displeasure we fear are mortal too. Their power over our circumstances will end, while we will still answer to God for the choices made under their influence.

Courage also requires prudence. We remain responsible for dependents, for careful speech, and for discerning the duties of our state. The courage we seek enables us to do the good required of us when doing it is costly. A parent may endure exhausting work for a child's sake; a physician may remain faithful to a vulnerable patient; a worker may refuse the dishonest act expected of everyone. Each must consider the actual duty and seek sound counsel where needed. We can ask Christ for the wisdom to recognize that duty and the strength to fulfill it despite the risks.

Hebrews also explains how Christ helps us in weakness. The Son became like his brothers so that he might be a merciful and faithful high priest; having suffered when tempted, he can help those undergoing trial (Heb 2:17--18). We therefore approach one who knows human vulnerability from within his sinless human life. The letter later invites us to draw near with confidence to the throne of grace, where mercy and timely help are given (Heb 4:14--16). We can bring our fear to him as soon as it arises and ask for strength to fulfill the duty before us.

Peter's failure shows how much we need such help. Before the Passion is complete, the disciple who followed Jesus into the high priest's courtyard denies knowing him. When the Lord turns and looks at him, Peter remembers and weeps (Lk 22:54--62). His collapse reveals how quickly a sincere attachment can yield under pressure. At the Lord's look he remembers the warning he had received and recognizes what he has done. When we discover cowardice in ourselves, we too can acknowledge the deed and turn to Christ for mercy. Remembering Peter can help us seek that mercy promptly instead of dwelling on how brave we once expected to be.

Remembering Peter can also help us examine how we respond to fear now. We cannot establish today exactly how we will feel under the pressure of a final illness. We can learn where we habitually seek refuge when frightened. We may grasp for control, withdraw, accuse someone, or make promises merely to escape discomfort. Bringing these habits into prayer and confession allows us to ask for help with them. Over time we can learn to pray when fear first arises, acknowledge what frightens us, and consider what faithfulness requires before acting to escape the discomfort.

Paul confesses the security of Christ's love amid suffering at the close of Romans~8. He names tribulation, distress, hunger, danger, persecution, and the sword; he even recalls the condition of those treated as sheep for slaughter. Then he confesses that neither death nor life nor any created power can separate us from the love of God in Christ Jesus (Rom 8:31--39). His confidence rests upon the Father who has given his Son, and upon Christ who died, rose, and intercedes. We can ask the interceding Christ to sustain our fidelity through the trials Paul names, including those we fear most.

We may therefore ask for the grace to obey Christ while still afraid. The annual remembrance of death gives us time to examine the occasions when we have preferred security to fidelity. We can acknowledge a particular failure, seek pardon, and ask how to fulfill the duty we avoided. We can also become more patient with someone whose courage is hard won. The Lord who knows our weakness gives us reason to persevere in this work without pretending that we have already overcome every fear.

